% =====================================================================
%  c03：原印刷页 22–30（PDF p.23–31）
%  内容：§2.2.2 末尾「Remarks」下编号列表的第 1–2 项（列表由 c02 页末的标题引
%        出，故此处另起 enumerate）、§2.2.3、§2.2.4、§2.2.5、§2.2.6（至页末）。
%  本片段起点没有编号标题（承接 §2.2.2 正文），故不补标题；四级标题均按编号
%  点数取 \subsubsection；原文中以斜体下划线标出的无编号小标题用 \paragraph 复刻。
%  本片段无插图；编号公式 (2.50)–(2.93) 连续。
% =====================================================================

\begin{enumerate}
  \item 一个关键之处在于：\eqref{eq:2.51} 与 \eqref{eq:2.52} 里那个无量纲的耦合常数 $k$ 是量子化的。人们有时会声称，在阿贝尔规范群的 Chern-Simons 理论中情形并非如此。若取 $G=\mathbb{R}$，或者更确切地说，只考虑那些相对平凡的、单连通的 3-流形（例如 $M_3=\mathbb{R}^3$），情形确实如此。但若规范群是 $U(1)$，而我们考虑拓扑上有意思的情形，那么它事实上必定是量子化的。

% [Tutor Note: 记号易错点：本列表第 1 项把耦合常数写作 $k$，但 \eqref{eq:2.51}、\eqref{eq:2.52} 印出的符号是 $\kappa$，全篇阿贝尔情形亦一律用 $\kappa$（$k$ 留给非阿贝尔情形，见 \eqref{eq:2.80} 的 $k=-2\kappa$）。引用「$k$ 是量子化的」这一结论时，应把它理解为 $\kappa\in\mathbb{Z}$（半整数情形见下一条），不要与第 3 章的能级 $k$ 混淆。]

  \item 这里我们应当提到，这个量子化的故事还有一个重要的精化。在物理的某些用途中（例如在分数量子 Hall 效应中），我们需要取 $\kappa$ 为半整数。如果我们不加思索地取 $\kappa=1/2$，那么
\begin{equation}
  \exp\left[\mathrm{i}\frac{1}{4\pi}\int_{M_3}AF\right]
  \tag{2.50}
  \label{eq:2.50}
\end{equation}

就不是规范不变的：这或许是最简单的例子——取 $M_3=T^3$，$F/2\pi=d\sigma^{1}d\sigma^{2}$，并作一个参数为 $\omega=2\pi d\sigma^{3}$ 的规范变换。这里 $\sigma^{i}\sim\sigma^{i}+1$ 是环面上的坐标。尽管如此，只要在定义理论时补充一些数据，就可以用规范不变的方式定义像 \eqref{eq:2.50} 这样的表达式。具体地说，若在 $M_3$ 上选定一个「旋量结构」（spin structure），就有可能定义路径积分中规范不变的测度。所付的代价是：作用量乃至整个理论随后会依赖于旋量结构的选取——尽管这种依赖是可计算的。我们将在下文第 2.5 节详细回到这个重要问题。
\end{enumerate}

% [Tutor Note: \eqref{eq:2.50} 的「不规范性」可以手算验证，这是理解旋量结构为何必要的最快途径：取大规范变换参数 $\omega=2\pi d\sigma^{3}$（它在 $S^1$ 方向绕一圈给出 $2\pi$ 的周期），代入 $\int_{M_3}\omega\wedge F=2\pi\cdot2\pi\int_{T^3}d\sigma^{3}\wedge d\sigma^{1}\wedge d\sigma^{2}=4\pi^{2}$，于是指数上的相位移动为 $\mathrm{i}\cdot\frac{1}{4\pi}\cdot4\pi^{2}=\mathrm{i}\pi$，即 $e^{\mathrm{i}\pi}=-1$。注意此处系数是 $1/4\pi$，而 \eqref{eq:2.51} 的 Chern-Simons 项系数是 $\kappa/2\pi$；两式相差因子 2 的归一化约定，换算时务必先统一。]

\subsubsection{三维 $U(1)$ 的 Maxwell-Chern-Simons 理论} % 原文: 2.2.3 3D Maxwell-Chern-Simons For $U(1)$

\paragraph{场与作用量} % 原文: 无编号小标题 Fields And Action（原以斜体下划线标示，非编号标题）

我们先从三维的「Maxwell-Chern-Simons」$U(1)$ 规范理论讲起。设 $M_3$ 是一个带定向的 Riemann 3-流形。于是作用量对路径积分的贡献通过因子 $e^{\mathrm{i}S}$ 给出，其中

\begin{equation}
  S=\int_{M_3}\left[-\frac{1}{8\pi^{2}e^{2}}dA\wedge\ast dA+\frac{\kappa}{2\pi}A\wedge dA\right]
  \tag{2.51}
  \label{eq:2.51}
\end{equation}

这里我们在 Lorentz 号差 $(-,+,+)$ 的 3-流形上工作，而 Lorentz 度规诱导出 Hodge 星 $\ast$。我们选取局部坐标，使得（视为整体非零 3-形式的）定向为 $\mathrm{vol}=dx^{0}\wedge dx^{1}\wedge dx^{2}$。

在欧氏理论中，作用量对路径积分的贡献是因子 $e^{-S}$，其中

\begin{equation}
  S_E=\int_{M_3}\left[-\frac{1}{8\pi^{2}e^{2}}dA\wedge\ast dA-\mathrm{i}\frac{\kappa}{2\pi}A\wedge dA\right]
  \tag{2.52}
  \label{eq:2.52}
\end{equation}

\textbf{Remark（评注）：}注意，与标准的 Maxwell 项不同，Chern-Simons 项要求选定一个定向。反转定向的微分同胚（例如宇称反转）会改变 Chern-Simons 作用量的符号。此外要注意：无论在欧氏号差还是 Minkowski 号差下，Chern-Simons 项对路径积分测度的贡献都是一个相位。

\paragraph{平面波与局域自由度} % 原文: 无编号小标题 Plane Waves And Local Degrees Of Freedom

Lorentz 号差下的运动方程是

\begin{equation}
  d\ast F-mF=0
  \tag{2.53}
  \label{eq:2.53}
\end{equation}

其中我们定义 $m:=4\pi ke^{2}$，它具有质量的量纲。考虑 Minkowski 空间 $M_3=\mathbb{M}^{1,2}$ 中的理论。为分析局域自由度，我们考虑平面波解：

\begin{equation}
  A=\epsilon_i(\vec p)e^{\mathrm{i}p\cdot x}+\text{c.c.}
  \tag{2.54}
  \label{eq:2.54}
\end{equation}

其中 $p\cdot x=p_0t+p_1x^{1}+p_2x^{2}$。于是 \eqref{eq:2.53} 变为：

\begin{equation}
  \begin{aligned}
    (p_1p_2+\mathrm{i}mp_0)\epsilon_1+(p_0^{2}-p_1^{2})\epsilon_2 &=0\\
    (p_0^{2}-p_2^{2})\epsilon_1+(p_1p_2-\mathrm{i}mp_0)\epsilon_2 &=0\\
    (p_1p_0+\mathrm{i}mp_2)\epsilon_1+(p_0p_2-\mathrm{i}mp_1)\epsilon_2 &=0
  \end{aligned}
  \tag{2.55}
  \label{eq:2.55}
\end{equation}

当且仅当

\begin{equation}
  p_0^{2}\left(p_0^{2}-p_1^{2}-p_2^{2}-m^{2}\right)=0
  \tag{2.56}
  \label{eq:2.56}
\end{equation}

时，$\epsilon_i$ 才存在非零解。解有两支。

\begin{enumerate}
  \item 若 $p_0\neq0$，则定义 $\omega(\vec p):=+\sqrt{\vec p^{2}+m^{2}}$。当 $p_0=\pm\omega(\vec p)$ 时规范被完全固定，而 $\epsilon_i$ 被确定到相差一个整体相乘的常数：
\begin{equation}
  \begin{aligned}
    \epsilon_1^{\pm}(\vec p) &=mp_1+\mathrm{i}p_0p_2\\
    \epsilon_2^{\pm}(\vec p) &=mp_2-\mathrm{i}p_0p_1
  \end{aligned}
  \tag{2.57}
  \label{eq:2.57}
\end{equation}

注意 $\epsilon_i=\epsilon_i^{+}=(\epsilon_i^{-})^{*}$。由于这两种极化相互关联，我们拥有的只是一个质量为 $|m|$ 的标量场的自由度。这是 [15, 14] 的主要结果。它推广了三维 $U(1)$ 规范场等价于一个紧致标量场这一众所周知的结论。

{\small\itshape ［原文标记：我是否应当解释 $U(1)$ 规范场与紧致标量场之间的对偶？］}% [原文待办标记，原文 p.23 页边竖排]

% [Tutor Note: \eqref{eq:2.55}–\eqref{eq:2.56} 的推导值得自己过一遍：三式中只有两个未知量 $\epsilon_1,\epsilon_2$，可解性要求该 $3\times2$ 系数矩阵的秩小于 2，即所有 $2\times2$ 子式同时为零，其公共零点正是 \eqref{eq:2.56} 的 $p_0^{2}\left(p_0^{2}-\vec p^{2}-m^{2}\right)=0$——其中 $p_0=0$ 一支对应后文第二支（平坦场），$p_0^{2}=\vec p^{2}+m^{2}$ 一支给出 \eqref{eq:2.57}。容易出错的是 \eqref{eq:2.55} 中 $\mathrm{i}mp_0$、$\mathrm{i}mp_1$、$\mathrm{i}mp_2$ 的来源：它们来自 $m\ast F$ 用 $\epsilon_{\mu\nu\lambda}$ 缩并时的符号与 $\mathrm{i}$ 的归属；另外原文的 $\epsilon_i=\epsilon_i^{+}=(\epsilon_i^{-})^{*}$ 只能读成「两支极化互为复共轭」，$i=1,2$ 仍是分量指标，不是自由参数。]

  \item 运动方程的另一支解是 $p_0=0$。此时 $p_1\epsilon_2-p_2\epsilon_1=0$，即 $F=0$，我们得到的是平坦规范场。这给出理论的拓扑扇区。在拓扑非平凡的情形中，我们可以用非局域算符、例如 Wilson 线，在 $\mathbb{M}^{1,2}$ 上探测这个扇区；在拓扑非平凡的其他物理情形中亦然。
\end{enumerate}

\paragraph{长距离/强耦合极限：拓扑扇区} % 原文: 无编号小标题 Long-Distance/Strong-Coupling Limit: The Topological Sector

耦合常数 $e^{2}$ 具有质量的量纲。在三维度规的共形重标度 $g_{\mu\nu}\to\Omega^{2}g_{\mu\nu}$ 下，作用量中第一项按 $\Omega^{-1}$ 标度，而第二项不变。注意两点：

% [原文注：原文此处写 "Note two things"，实际列出三项；第 1 项在原 p.23 页末，第 2、3 项在原 p.24 页首。译文照原文保留这一不一致。]

\begin{enumerate}
  \item 长距离极限是这样得到的：在 $M_3$ 上取固定的度规 $g_{\mu\nu}$，并令 $\Omega\to\infty$。

  \item 在该极限下，至少形式上，第一项消失而只剩第二项：我们得到一个 Chern-Simons 理论。

  \item 把 $g_{\mu\nu}\to\Omega^{2}g_{\mu\nu}$ 与把 $e^{2}\to\Omega e^{2}$ 是等价的。所以长距离极限与强耦合极限是同一件事。
\end{enumerate}

当 $e^{2}\to\infty$ 时质量 $m\to\infty$，传播子趋于零，从而杀死所有 Feynman 图与所有局域关联。所以像

\begin{equation}
  \left\langle F_{\mu_1\nu_1}(p_1)F_{\mu_2\nu_2}(p_2)\cdots F_{\mu_n\nu_n}(p_n)\right\rangle
  \tag{2.58}
  \label{eq:2.58}
\end{equation}

这样的局域规范不变可观测量的关联函数将全部为零。换一种想法来看：当插入点 $p_1,\dots,p_n$ 彼此相距很远时，插入一组完备的态，并利用这些态是一个大质量标量场的态这一事实，可以看出该关联函数必须按

\begin{equation}
  \operatorname*{Max}_{i\neq j}e^{-\mathrm{const.}\,|m|\cdot\|p_i-p_j\|}
  \tag{2.59}
  \label{eq:2.59}
\end{equation}

的方式衰减。

但这并不使该极限下的理论变得平凡。设 $R$ 是规范群 $U(1)$ 的一个有限维表示，而 $\gamma\subset M_3$ 是一条闭的定向圈，则我们可以定义经典可观测量

\begin{equation}
  W(R,\gamma):=\operatorname{Tr}_R\exp\left(\mathrm{i}\oint_{\gamma}A\right)
  \tag{2.60}
  \label{eq:2.60}
\end{equation}

当然，对 $U(1)$ 而言每个有限维表示都完全可约为一维不可约表示之和，而不可约表示为 $\rho_n(z)=z^{n}$（$n\in\mathbb{Z}$），故不失一般性，我们可以只考虑经典 Wilson 线

\begin{equation}
  W(n,\gamma):=\exp\left(\mathrm{i}n\oint_{\gamma}A\right)
  \tag{2.61}
  \label{eq:2.61}
\end{equation}

{\small\itshape ［原文标记：$R$ 这个记号选得不好，因为它看起来像个半径。］}% [原文待办标记，原文 p.24 页边竖排，原文小写作 r]

在取 $e^{2}\to\infty$ 之前，这些 Wilson 线在路径积分中的关联函数会强烈依赖于时空中的曲线 $\gamma$，也会依赖于 $\gamma$ 上各点与 $F_{\mu\nu}(p)$ 这类局域算符的邻近程度。然而 \eqref{eq:2.61} 的关联函数在 $e^{2}\to\infty$ 极限下将只依赖于 $\gamma$ 到「同痕」（isotopy）为止。这就是说，它们只依赖于 $\gamma$ 在 $M_3$ 中（挖去其他算符所在位置之后）的连续形变类。特别地，你不能把一条曲线 $\gamma$ 穿过另一条曲线 $\gamma'$ 而作形变。为了推出这些结论，我们现在转而对「纯」Chern-Simons 项作量子化，即在作用量 \eqref{eq:2.52} 中形式上取极限 $e^{2}\to\infty$。

\subsubsection{$U(1)$ Chern-Simons 理论的形式路径积分} % 原文: 2.2.4 The Formal Path Integral Of The $U(1)$ Chern-Simons Theory

我们现在考虑由 Chern-Simons 项自身定义路径积分测度的 Chern-Simons 理论。也就是说，我们想要赋予下述路径积分以意义：

\begin{equation}
  \int_{\mathcal{A}/\mathcal{G}}[dA]\,e^{2\pi\mathrm{i}\frac{\kappa}{4\pi^{2}}\int_{M_3}A\wedge dA}\prod_{\alpha}W(n_{\alpha},\gamma_{\alpha})
  \tag{2.62}
  \label{eq:2.62}
\end{equation}

其中积分形式上遍历固定 $M_3$ 上线丛的联络的规范等价类所构成的空间 $\mathcal{A}/\mathcal{G}$。

几点直接评注：

% [Tutor Note: 归一化核对：\eqref{eq:2.62} 指数上的系数是 $2\pi\mathrm{i}\cdot\frac{\kappa}{4\pi^{2}}=\mathrm{i}\frac{\kappa}{2\pi}$，恰好等于 \eqref{eq:2.51} 中 Chern-Simons 项 $\frac{\kappa}{2\pi}A\wedge dA$ 的系数，所以 \eqref{eq:2.62} 就是 $e^{\mathrm{i}S_{CS}}$。若把 \eqref{eq:2.62} 误读为 $e^{2\pi\mathrm{i}\kappa\int A\wedge dA}$（丢掉 $4\pi^{2}$），$\kappa$ 的量子化条件就会与第 2.2.2 节 (2.48)–(2.49) 的整性论证相差 $4\pi^{2}$ 的因子。另外，正因为系数里没有 $\hbar$，$\kappa$ 在此充当 $1/\hbar$（见下面第 2 条），半经典展开的小参数是 $1/\kappa$。]

\begin{enumerate}
  \item 这是一个具有二次作用量的自由场论。路径积分可以被严格定义，一切都应当完全可解。没有任何借口。

  \item 耦合常数 $\kappa$ 扮演 $1/\hbar$ 的角色。

  \item 记下与源 $J$ 耦合时的经典运动方程是有用的。若把 $J$ 视为一个 2-形式，则
\begin{equation}
  \frac{\kappa}{2\pi}\int AdA+\int A\wedge J
  \tag{2.63}
  \label{eq:2.63}
\end{equation}

给出运动方程

\begin{equation}
  F=-\frac{\pi}{\kappa}J
  \tag{2.64}
  \label{eq:2.64}
\end{equation}

  \item 作用量 $\int_{M_3}A\wedge dA$ 不涉及任何时空度规。这提示该理论是一个「拓扑场论」，其中所有跃迁函数、振幅、关联函数等等都与距离无关。（稍后我们会给出一个更形式化、更数学化的定义。）

  \item 不过，不应仓促地跳到这个结论，因为毕竟还需要定义测度，而在定义算符的关联函数时可能需要某种正规化。确实，在量子层面会存在某种非常温和的度规依赖，但它很容易被描述并隔离出来，本质上这将证明是一个真正的拓扑场论。参见下文第 2.2.17 节与第 3.2.1 节。

  \item 作为我们将遇到的量子微妙之处的一个例子：Wilson 线可观测量 $W(n,\gamma)$ 的定义会依赖于被称为「标架」（framing）的额外数据；事实上它们是与带子（ribbon）、而非与曲线相联系的观测量。参见第 2.2.16 节。
\end{enumerate}

\subsubsection{通往态 Hilbert 空间的最初几步} % 原文: 2.2.5 First Steps To The Hilbert Space Of States

接下来若干节（直到下文第 2.2.14 节之前，不含该节）将处理如下问题：设 $M_3=\Sigma_2\times\mathbb{R}$，其中 $\Sigma_2$ 是一个闭的拓扑曲面。应当存在一个与 $\Sigma_2$ 相联系的（Hilbert）量子态空间。由于理论依赖于 $M_3$ 的定向，而这样的定向由 $\Sigma_2$ 的一个定向 $o(\Sigma_2)$ 所决定，我们预期该空间依赖于 $\Sigma_2$、$o(\Sigma_2)$，当然还有 $\kappa$。我们将把它记为

\begin{equation}
  \mathcal{H}(\Sigma_2,o(\Sigma_2),\kappa)
  \tag{2.65}
  \label{eq:2.65}
\end{equation}

\textbf{Remark（评注）：}稍后，当我们在 $\Sigma_2$ 上引入复结构 $J$ 时，可以把它写成 $\mathcal{H}(\Sigma_2,J,\kappa)$，而它将是 $\Sigma_2$ 上复结构模空间之上的一个 Hilbert 空间丛的纤维。

关于作用量

\begin{equation}
  \int_{M_3}A\wedge dA=\int_{M_3}d^{3}x\,\epsilon^{\mu\nu\lambda}A_{\mu}\partial_{\nu}A_{\lambda}
  \tag{2.66}
  \label{eq:2.66}
\end{equation}

我们应当作出的第一个观察是：它关于「时间」是一阶的。若我们选取局部坐标 $(x^{0},x^{1},x^{2})$ 并以 $x^{0}$ 充当时间，则可写 $A=A_0dx^{0}+A_s$，其中 $A_s$ 只含「空间分量」，并且

\begin{equation}
  A\wedge dA=A_sdx^{0}\partial_0A_s-A_0F_s
  \tag{2.67}
  \label{eq:2.67}
\end{equation}

其中 $F_s$ 是空间分量。

有鉴于此，我们将专注于 Hamilton 量子化——即通过选定时间来破坏 Lorentz 对称性。于是我们考虑 3-流形取如下形式

\begin{equation}
  M_3=\Sigma_2\times\mathbb{R}
  \tag{2.68}
  \label{eq:2.68}
\end{equation}

其中 $\Sigma_2$ 是一个二维定向曲面。我们可以写出

{\small\itshape ［原文标记：记号问题——$\Sigma_2$ 表示定向曲面；后文对带有复结构的曲面又用 $(\Sigma_2,J)$，即 Riemann 曲面？？］}% [原文待办标记，原文 p.26 页边竖排]

\begin{equation}
  \int_{M_3}A\wedge dA=-\int dx^{0}\left(\int_{\Sigma_2}A_s\partial_0A_s\right)-\int dx^{0}\left(\int_{\Sigma_2}A_0F_s\right)
  \tag{2.69}
  \label{eq:2.69}
\end{equation}

% [Tutor Note: 从 \eqref{eq:2.67} 到 \eqref{eq:2.69} 出现了整体负号，且 $A_0$ 项与 $A_s\partial_0A_s$ 项的符号必须一起核对。检查办法：固定 $\mathrm{vol}=dx^{0}\wedge dx^{1}\wedge dx^{2}$，把 $A=A_0dx^{0}+A_s$ 直接代入 $A\wedge dA$ 展开（注意 $dx^{0}$ 与空间 1-形式交换次序会带出负号），再写成 $\int_{M_3}=\int dx^{0}\int_{\Sigma_2}$。这一步的符号会一路传到 \eqref{eq:2.89}–\eqref{eq:2.92} 的动量映射上，作者本人在 \eqref{eq:2.91} 旁也自注 NEED TO CHECK SIGNS，改符号时必须整链一起改。]

现在，作用量关于时间导数是一阶的。这相当重要。一般地，关于时间导数为一阶的作用量形如

\begin{equation}
  S=\int dt\,\lambda_i(\phi)\frac{d\phi^{i}}{dt}
  \tag{2.70}
  \label{eq:2.70}
\end{equation}

我们应当把该作用量看作相空间 $\mathcal{P}$ 中轨迹的作用量。相空间就是辛流形。关于辛几何的一些定义与事实的综述见附录 A。注意作用量的变分为

\begin{equation}
  \delta S=\int dt\left(\partial_i\lambda_j-\partial_j\lambda_i\right)\delta\phi^{i}\frac{d\phi^{j}}{dt}
  \tag{2.71}
  \label{eq:2.71}
\end{equation}

而若作用量是非退化的，我们就可以读出辛形式

\begin{equation}
  \omega=\left(\partial_i\lambda_j-\partial_j\lambda_i\right)\delta\phi^{i}\wedge\delta\phi^{j}
  \tag{2.72}
  \label{eq:2.72}
\end{equation}

换言之，若（局部地）写 $\omega=d\lambda$（$\lambda$ 为某个 1-形式），则沿相空间中路径 $\gamma$ 的轨迹的作用量为

\begin{equation}
  S=\int_{\gamma}\lambda-\int dt\,H(\gamma(t))
  \tag{2.73}
  \label{eq:2.73}
\end{equation}

局部地，在 Darboux 坐标下可写 $\lambda=p_idq^{i}$，于是我们的作用量为

\begin{equation}
  S=\int p_i\frac{dq^{i}}{dt}dt-\int H\,dt=\int p\,\dot q\,dt-\int H(p(t),q(t))\,dt
  \tag{2.74}
  \label{eq:2.74}
\end{equation}

我们由此得出结论：

\begin{quote}
  \itshape 我们应当把 $\Sigma$ 上的规范场（即 $P$ 上的联络）所构成的空间 $\mathcal{A}(P\to\Sigma)$ 看作相空间，而作用量就是一个在相空间上运动、并受约束的粒子的作用量。
\end{quote}

% [Tutor Note: \eqref{eq:2.72} 对 $i,j$ 作全求和，比起 $\omega=\tfrac12\omega_{ij}dx^{i}\wedge dx^{j}$ 的约定来说少了一个 $1/2$——同一个 $\omega$ 被数了两遍。这不影响此后物理，但在把 \eqref{eq:2.79} 反解成 Poisson 括号时必须把约定定准。最快的自检是走一遍 \eqref{eq:2.79}$\to$\eqref{eq:2.85}：由 $\delta A_s\wedge\delta A_s=2\,\delta A_1\delta A_2\,dx^{1}\wedge dx^{2}$ 得 $\omega=\frac{\kappa}{\pi}\int\delta A_1\delta A_2$，故按 $\omega=\tfrac12\omega^{ij}\delta A_i\delta A_j$ 读出 $\omega^{12}=\kappa/\pi$，其逆给出 $\{A_1,A_2\}=\pm\frac{\pi}{\kappa}\delta^{(2)}(x,y)$——绝对值正是 \eqref{eq:2.85} 的 $\pi/\kappa$。若你的符号与之相反，说明你用的 $\{f,g\}=\omega^{ij}\partial_if\partial_jg$ 与作者的约定整体差一个符号。]

\subsubsection{关于相空间量子化与 Hamilton 约化的一般评注} % 原文: 2.2.6 General Remarks On Quantization Of Phase Space And Hamiltonian Reduction

\paragraph{「量子化一个相空间」是什么意思？} % 原文: 无编号小标题 What Does It Mean To "Quantize A Phase Space?"

设给定了辛流形 $\mathcal{P}$。也就是说，它是一个配有闭 2-形式 $\omega=\frac12\omega_{ij}(x)dx^{i}\wedge dx^{j}$ 的流形，使得反对称矩阵 $\omega_{ij}(x)$ 在流形上每一点 $x$ 处都可逆。我们想要「量子化 $\mathcal{P}$」。让我们精确说明这句话的意思：

首先，我们考虑一族辛形式 $\omega^{(k)}=k\omega$，以便讨论半经典极限。在我们的例子里 $k$ 将是整数。

我们想给每个 $k$ 指定一个 Hilbert 空间 $\mathcal{H}_k$，并对适当的函数 $f:\mathcal{P}\to\mathbb{C}$ 指定 $\mathcal{H}_k$ 上相应的线性算符 $Q^{(k)}(f)$。我们要求：

\begin{equation}
  \lim_{k\to\infty}\left\|Q^{(k)}(f)Q^{(k)}(g)-Q^{(k)}(fg)\right\|<\infty
  \tag{2.75}
  \label{eq:2.75}
\end{equation}

以及

\begin{equation}
  \left[Q^{(k)}(f),Q^{(k)}(g)\right]=-\frac{\mathrm{i}}{k}Q^{(k)}(\{f,g\})+\mathcal{O}(1/k^{2})
  \tag{2.76}
  \label{eq:2.76}
\end{equation}

其中

\begin{equation}
  \{f,g\}=\omega^{ij}\partial_if\partial_jg
  \tag{2.77}
  \label{eq:2.77}
\end{equation}

是 Poisson 括号。此外我们还要求

\begin{equation}
  \left(Q^{(k)}(f)\right)^{\dagger}=Q^{(k)}(f^{*})
  \tag{2.78}
  \label{eq:2.78}
\end{equation}

从而实函数对应于自伴算符。

\paragraph{如何量子化相空间} % 原文: 无编号小标题 How To Quantize Of Phase Space

量子化相空间有许多途径。据我们所知，并不存在已知的一般程序。\footnote{如果我们满足于在 $1/k$ 的形式幂级数中工作，那么 Fedosov 与 Kontsevich 的结果（尤其是 Cattaneo 与 Felder 对 Kontsevich 定理的物理解释）确实提供了答案。} 不过，若给定更多结构，则有一些一般方法：

\begin{enumerate}
  \item \emph{Schrödinger 量子化}：有时存在 $p_i,q^{i}$ 的整体分离。最简单的情形是 $\mathcal{P}=T^{*}M$，其中 $M$ 是流形。此时 $\mathcal{H}=L^{2}(M)$。态是归一化波函数 $\Psi(q^{i})$，而 $p_i=-\mathrm{i}\hbar\nabla_{q^{i}}$ 是一阶微分算符。在算符排序的意义上，我们可以把经典可观测量——即 $\mathcal{P}$ 上的函数——翻译为 $\mathcal{H}$ 上的算符。\footnote{在我们对 Schrödinger 量子化的应用中，会有一个重要的额外微妙之处：动量本身是周期的，因此我们将要量子化的并不是形如 $T^{*}M$ 的相空间。}

  \item 有时相空间上带有允许量子化的额外结构。一个很好的例子是：$\mathcal{P}$ 同时又是 Kähler 流形，且存在一个全纯线丛 $L\to\mathcal{P}$ 并配有适当的联络。此时我们可以使用\emph{几何量子化}、\emph{Kähler 量子化}或 \emph{Berezin-Toeplitz 量子化}的方法。附录 B 回顾了几何量子化与相干态理论的一些要点。
\end{enumerate}

这些途径并非普遍适用。由于四维流形理论的进展，现在已有大量紧致辛四维流形的例子，它们不容许 Kähler 流形的结构。尽管如此，在这些讲义中我们将遇到的辛空间的例子里，上述两种程序就够用了。事实上，我们主要将使用 Kähler 量子化方法。

\paragraph{Hamilton 约化} % 原文: 无编号小标题 Hamiltonian Reduction

关于辛空间在辛群作用下的 Hamilton 约化的背景，见附录 A。

当我们量子化一个带有生成辛 Lie 群作用的一阶约束的相空间时，有两种量子化方式：

\begin{enumerate}
  \item 先在经典层面作 Hamilton 约化：算出满足 $\mu=0$ 的场空间，再算出商 $\mu^{-1}(0)/G=\bar{\mathcal{P}}$，然后量子化这个商空间。总得想办法。

  \item 先在大相空间「楼上」量子化，然后用算符约束对允许的态施加约束。在我们的情形中，我们用 \eqref{eq:2.87} 量子化规范场的空间。然后把 Gauss 定律施加到波泛函 $\Psi[A]$ 上，即要求对规范群中所有 $g$ 都有 $U(g)\Psi=\Psi$。

  \item 也有混合情形：先对某些自由度在经典层面施加约束，然后量子化剩余的自由度，然后再量子化。\footnote{这适用于非阿贝尔理论的某些「阿贝尔化程序」。{\small\itshape ［原文标记：NEED TO CITE SOME REFERENCES FOR THIS——此处需补引文献；并包括 [17] 的最后一节。］}}
\end{enumerate}

先验地（\emph{a priori}）并不显然，不同程序会给出相同的结果。人们当然希望如此；所幸在 Chern-Simons 理论中，一切都顺利。

{\small\itshape ［原文标记：请给一个这些程序互不对易的例子？］}% [原文待办标记，原文 p.28 页边竖排]

\paragraph{应用到 Chern-Simons 理论} % 原文: 无编号小标题 Application To Chern-Simons Theory

该相空间上的辛形式为

\begin{equation}
  \omega=\frac{\kappa}{2\pi}\int_{\Sigma}\delta A_s\wedge\delta A_s
  \tag{2.79}
  \label{eq:2.79}
\end{equation}

我们在此记下向非阿贝尔理论的推广。使用第 3.1 节讨论的 Lie 代数 $\mathfrak{g}$ 上一个 Ad-不变迹的归一化，我们有

\begin{equation}
  \omega=\frac{k}{4\pi}\int_{\Sigma}\operatorname{tr}\left(\delta A_s\wedge\delta A_s\right)
  \tag{2.80}
  \label{eq:2.80}
\end{equation}

于是，从非阿贝尔情形过渡到 $U(1)$，只需丢掉 tr 并取 $k=-2\kappa$。要理解这个负号，回想：对 $G=U(1)$ 我们取 $A$ 为实值，而对一般的 $G$ 我们取 $A$ 取值于 Lie 代数。这就解释了 \eqref{eq:2.79} 与 (3.37) 之间的相对符号。

对你们中的许多人来说，记号 $\delta A_s$ 是不言自明的：它是空间上的一个 1-形式，同时又是全体规范场构成的空间上的一个 1-形式，对应于联络的一个小变分。

% [Tutor Note: 原文此处两处引用「(3.37)」（本页的「\eqref{eq:2.79} 与 (3.37) 之间的相对符号」以及原 p.29 的「So, (3.37) is shorthand for」）都是笔误，指的应是紧邻的 \eqref{eq:2.80} 与 \eqref{eq:2.84}：第 3.1 节的 (3.37) 是非阿贝尔作用量的归一化，与本处讨论的辛形式同源却并非同一式。按原文转写未作改动。又：由 \eqref{eq:2.80} 经 $k=-2\kappa$ 回到 \eqref{eq:2.79} 时，$\frac{k}{4\pi}=-\frac{\kappa}{2\pi}$，恰好与 \eqref{eq:2.79} 差一个整体负号——这个相对符号完全由「$A$ 在 $U(1)$ 情形取实值、在一般 $G$ 情形取值于 $\mathfrak{g}$」这一约定差异所造成，切勿当作笔误改掉。]

\textbf{Remark（评注）：}写给数学上爱较真的读者（\emph{For mathematically-minded fussbudgets}）。大多数物理学家看到 \eqref{eq:2.79} 这样的表达式，就知道该怎么做；但对于那些开始追问 $\delta A$ 这样的符号究竟是什么意思、又「住在哪里」的人，下面的评注或许有用。给定主丛 $P\to M$，$P$ 上全体联络构成的空间 $\mathcal{A}=\mathrm{Conn}(P\to M)$ 在每一点 $A\in\mathcal{A}$ 处的切空间典范地是

\begin{equation}
  T_A\mathcal{A}\cong\Omega^{1}(M;\operatorname{ad}P)
  \tag{2.81}
  \label{eq:2.81}
\end{equation}

联络空间是仿射空间，因此切空间「不依赖于」点 $A$ 是合理的。为说明这一点，注意若我们有一族连续联络 $d+A(t)$，则时刻 $t+\epsilon$ 与 $t$ 之差为

\begin{equation}
  (d+A(t+\epsilon))-(d+A(t))=A(t+\epsilon)-A(t)=\epsilon\alpha+\mathcal{O}(\epsilon^{2})
  \tag{2.82}
  \label{eq:2.82}
\end{equation}

当 $\epsilon\to0$，其中 $\alpha$ 是某个整体定义的、取值于伴随表示的 1-形式。现在设 $\gamma(t)=d+A(t)$ 是一条联络路径，使得切向量为 $\dot\gamma(0)=\alpha$。此时 $\delta A$ 就是满足下式的 1-形式：

\begin{equation}
  \delta A(\dot\gamma(0))=\alpha
  \tag{2.83}
  \label{eq:2.83}
\end{equation}

当你看到 $\mathcal{A}$ 上以 $\delta A$ 写出的微分形式时，若想要数学上精确的版本，就把它作用于由 $\alpha_1,\alpha_2,\dots$ 构成的多向量场，并使用 \eqref{eq:2.83}（再加上一些常识），你就会得到数学文献中见到的那类公式。

换一种想法是这样的：若我们设想把丛平凡化、选取局部坐标并选取 Lie 代数 $\mathfrak{g}$ 的一组基 $t_a$，则在 $\mathcal{A}$ 上有局部「坐标」$A^{a}_{\mu}(x)$，其中 $a,\mu,x$ 都是「指标」。如果你对 $x$ 这样的连续指标感到不适，就用格点规范理论把 $(x,\mu)$ 合成一个链变量。于是我们有局部定义的微分 $\delta A^{a}_{\mu}(x)$，其中 $\delta$ 表示联络空间上的外微分。所以 (3.37) 是下式的简写：

\begin{equation}
  \omega|_{\mathcal{A}}=\frac{\kappa}{2\pi}\int_{\Sigma}\delta A^{a}_{i}(x)\wedge\delta A^{a}_{j}(x)dx^{i}\wedge dx^{j}
  \tag{2.84}
  \label{eq:2.84}
\end{equation}

其中我们取 $t_a$ 为 $\mathfrak{g}$ 在非退化形式 tr 下的一组标准正交基。

把 \eqref{eq:2.79} 所隐含的 Poisson 括号用局部坐标写出，即为

\begin{equation}
  \{A_i(x),A_j(y)\}=\frac{\pi}{\kappa}\epsilon_{ij}\delta^{(2)}(x,y)
  \tag{2.85}
  \label{eq:2.85}
\end{equation}

这里 $\delta^{(2)}(x,y)$ 是权为一的 Dirac delta 函数，而我们需要在 $\Sigma_2$ 上选一个定向来定义 $\epsilon_{ij}$。我们选取局部坐标使得 $\epsilon_{12}=1$。要把归一化搞对，回想：若 $\omega=dp\wedge dq$，则 $\{p,q\}=1$，这会有所帮助。

非阿贝尔情形下我们有：

\begin{equation}
  \{A^{a}_{i}(x),A^{b}_{j}(y)\}=\frac{2\pi}{k}\epsilon_{ij}\delta^{ab}\delta^{(2)}(x,y)
  \tag{2.86}
  \label{eq:2.86}
\end{equation}

其中 $A=A^{a}_{i}t_adx^{i}$，而 $t^{a}$ 是关于迹 tr 的一组标准正交基（其中{\small\itshape ［原文标记：务必核实这在几何上是对的——右端应是正确丛的截面；另外 $\epsilon$ 既是局部规范参数，又是张量 $\epsilon_{ij}$。］}% [原文待办标记，原文 p.30 页边竖排；此行原文在「tr.」之后还残留一个孤立的斜体 where，系插入边注后遗留，未译出]
$x,y\in\Sigma_2$）。于是，在最直截了当的量子化方案中，我们会说

\begin{equation}
  \left[\hat A_i(x),\hat A_j(y)\right]=-\frac{\mathrm{i}\pi}{\kappa}\epsilon_{ij}\delta^{(2)}(x,y)
  \tag{2.87}
  \label{eq:2.87}
\end{equation}

此外，相空间上还有一个群作用：空间规范变换群。由 $\epsilon:M_3\to\mathbb{R}$ 给出的无穷小规范变换生成 $\delta A=D_A\epsilon=d\epsilon+[A,\epsilon]$。几何上，存在一个向量场 $V(\epsilon)$，它在 $\mathcal{A}$ 上的 1-形式 $\delta$ 处的取值为

\begin{equation}
  V(\epsilon)(\delta A)=D_A\epsilon
  \tag{2.88}
  \label{eq:2.88}
\end{equation}

利用辛形式 \eqref{eq:2.79} 可以算出相应的动量映射（moment map）：

\begin{equation}
  \begin{aligned}
    \iota(V(\epsilon))\omega&=\frac{k}{2\pi}\int_{\Sigma}\operatorname{Tr}D_A\epsilon\,\delta A\\
    &=\frac{k}{2\pi}\int_{\Sigma}\left[d\left(\operatorname{tr}\epsilon\,\delta A\right)-\operatorname{tr}\epsilon\,D_A\delta A\right]\\
    &=-\frac{k}{2\pi}\int_{\Sigma}\operatorname{tr}\left(\epsilon D_A\delta A\right)+\frac{k}{2\pi}\int_{\Sigma}d\left(\operatorname{tr}\epsilon\,\delta A\right)
  \end{aligned}
  \tag{2.89}
  \label{eq:2.89}
\end{equation}

在 $\partial\Sigma_2=\varnothing$ 的情形，我们有

\begin{equation}
  \iota(V(\epsilon))\omega=\delta\left(-\frac{k}{2\pi}\int_{\Sigma}\operatorname{Tr}\epsilon F\right)
  \tag{2.90}
  \label{eq:2.90}
\end{equation}

所以：

{\small\itshape ［原文标记：需要核对符号！］}% [原文待办标记，原文 p.30 页边竖排]

\begin{equation}
  \left\langle\mu,\epsilon\right\rangle=-\frac{k}{2\pi}\int_{\Sigma_2}\operatorname{tr}\left(\epsilon F\right)
  \tag{2.91}
  \label{eq:2.91}
\end{equation}

也就是说，对 $\Sigma_2$ 上所有的点 $p$，

\begin{equation}
  \{\mu(\epsilon),A_i(p)\}=(-D_A\epsilon)_i(p)
  \tag{2.92}
  \label{eq:2.92}
\end{equation}

这里我们把它写成一种可推广到非阿贝尔情形的方式。

% [Tutor Note: \eqref{eq:2.89}–\eqref{eq:2.91} 有三处极易出错。其一，$\operatorname{Tr}$ 与 $\operatorname{tr}$ 在原文中混用：\eqref{eq:2.89} 首行与 \eqref{eq:2.90} 用 $\operatorname{Tr}$，\eqref{eq:2.89} 后两行与 \eqref{eq:2.91} 用 $\operatorname{tr}$——只有取同一归一化的迹，才能把 \eqref{eq:2.89} 的首行与末行、以及 \eqref{eq:2.90} 与 \eqref{eq:2.91} 对齐。其二，\eqref{eq:2.89} 第二行到第三行用到了 $d(\operatorname{tr}\epsilon\delta A)$ 是全导数、在 $\partial\Sigma_2=\varnothing$ 时积分为零，这正是 \eqref{eq:2.90} 成立的前提；若边界非空就必须保留该项（第 3.1.1 节正是处理这一情形）。其三，作者在 \eqref{eq:2.91} 旁自注 NEED TO CHECK SIGNS：\eqref{eq:2.91} 的整体符号与 \eqref{eq:2.92} 的 $\{\mu(\epsilon),A_i\}=(-D_A\epsilon)_i$ 是绑定的，改一处必须同时改另一处。]

现在回想，在 Hamilton 约化中我们考虑把 $\mu=0$ 的子空间对对称群取商。在我们的情形中，$\mu=0$ 意味着联络是平坦的，$F=0$，因此 Hamilton 约化给出的相空间是

\begin{equation}
  \mathcal{M}_{flat}=\{A\in\mathcal{A}\,|\,F(A)=0\}/\mathcal{G}
  \tag{2.93}
  \label{eq:2.93}
\end{equation}

于是我们得到一个关键结果：\emph{相空间的 Hamilton 约化就是 $\Sigma_2$ 上\underline{平坦}规范场的模空间。}

\textbf{Remarks（注记）：}

% 原 p.30 页末的「Remarks」编号列表在此被切断：第 2 项起接续于下一片段（c04）。
\begin{enumerate}
  \item 在简单的 $U(1)$ 情形，所有关于 Hamilton 约化的这番谈论有点杀鸡用牛刀：只需注意到 $A_0$ 在作用量中充当 Lagrange 乘子，因此路径积分中对 $A_0$ 的积分给出一个 delta 函数 $\delta(F_{12})$，把作用量局域到平坦规范场上。
\end{enumerate}
